% =====================================================================
% IV-C: kết quả cấp độ 2. Các dòng ablation (cạnh xáo trộn, nhiễu Gauss, β > 0) ở Bảng tab:l2-ablation, Mục III-B.
\subsection{Kết quả cấp độ 2: MKG-W và NativE}
\label{sec:l2}

\textbf{Bước a: tái lập link prediction của NativE.}
Nhóm chạy lại NativE trên MKG-W gốc với đúng mã nguồn và siêu tham số của tác giả. Tiêu chí đặt trước khi chạy: chênh không
quá khoảng 1 điểm thì coi là tái lập được. Kết quả lệch +0,41 MRR và +0,61 Hits@1 (Bảng~\ref{tab:l2-repro}), nên bước a đạt.
Bản \texttt{rmask}, đã che đặc trưng ảnh và chữ của 1\,784 thực thể val/test retrieval, cho cùng mức MRR, tức việc che
không làm giảm chất lượng link prediction. Mỗi lần huấn luyện mất khoảng 1,1 giờ trên một GPU RTX 4090. Tham số
\texttt{-seed} của mã nguồn chỉ đặt seed cho PyTorch; bộ lấy mẫu âm viết bằng C không được seed, nên chuỗi mẫu âm giống nhau
ở mọi seed. Độ lệch chuẩn qua 3 seed vì vậy chỉ phản ánh khởi tạo và nhiễu của phần đối kháng, không bao gồm biến thiên do
lấy mẫu âm.

\begin{table}[t]
\centering
\caption{Link prediction trên tập test KGC của MKG-W (filtered). Bài báo chỉ công bố MRR và Hits@1 cho MKG-W. Dòng rmask:
3 seed.}
\label{tab:l2-repro}
\footnotesize
\setlength{\tabcolsep}{3pt}
\resizebox{\columnwidth}{!}{%
\begin{tabular}{lcccc}
\toprule
 & MRR & Hits@1 & Hits@3 & Hits@10 \\
\midrule
Bài báo NativE & 36,58 & 29,56 & --- & --- \\
Tái lập, MKG-W gốc (seed 42) & 36,99 & 30,17 & 39,95 & 49,32 \\
MKG-W-rmask (seed 0, 1, 2) & \pmstd{36,99}{0,24} & \pmstd{29,98}{0,11} & \pmstd{40,30}{0,33} & \pmstd{49,73}{0,54} \\
\bottomrule
\end{tabular}}
\end{table}

\begin{table}[t]
\centering
\caption{Kết quả truy vấn trên MKG-W (892 truy vấn mỗi split, pool 892 ảnh). $\alpha,\beta$ chọn trên val; test chấm một lần.
Hai dòng NativE dùng riêng tín hiệu của NativE ($\beta = 1$).}
\label{tab:l2-results}
\footnotesize
\setlength{\tabcolsep}{2.5pt}
\resizebox{\columnwidth}{!}{%
\begin{tabular}{llcccc}
\toprule
 & Dòng & R@1 & R@5 & R@10 & MRR \\
\midrule
\multirow{4}{*}{\rotatebox{90}{Val}} & CLIP thuần & 24,33 & 44,39 & 53,81 & 34,22 \\
 & + đếm cạnh (dòng chính) & \textbf{32,40} & \textbf{50,78} & \textbf{57,62} & \textbf{41,18} \\
 & + điểm NativE & \pmstd{28,25}{0,11} & \pmstd{45,85}{0,51} & \pmstd{54,37}{0,70} & \pmstd{36,97}{0,08} \\
 & + cosine NativE & \pmstd{26,64}{0,26} & \pmstd{45,96}{0,39} & \pmstd{54,30}{0,62} & \pmstd{35,98}{0,15} \\
\midrule
\multirow{4}{*}{\rotatebox{90}{Test}} & CLIP thuần & 24,89 & 44,73 & 54,48 & 34,61 \\
 & + đếm cạnh (dòng chính) & \textbf{29,82} & \textbf{51,46} & \textbf{59,42} & \textbf{39,89} \\
 & + điểm NativE & \pmstd{27,24}{0,19} & \pmstd{46,52}{0,11} & \pmstd{56,28}{0,11} & \pmstd{36,67}{0,14} \\
 & + cosine NativE & \pmstd{26,08}{0,45} & \pmstd{44,92}{0,53} & \pmstd{54,22}{0,95} & \pmstd{35,46}{0,58} \\
\bottomrule
\end{tabular}}
\end{table}

\textbf{Bước b: kết quả truy vấn} (Bảng~\ref{tab:l2-results}). ``Đếm cạnh'' là kênh \texttt{triple}; ``điểm NativE'' là kênh
\texttt{kgc} (dòng \texttt{kgc\_only}); ``cosine NativE'' là kênh \texttt{emb} (dòng \texttt{native\_emb}). Khi kết hợp
NativE với kênh đếm cạnh, val chọn $\beta = 0$, nên các dòng kết hợp trùng với dòng chính (Bảng~\ref{tab:l2-ablation}). Bảng
này đủ ba dòng của khung so sánh: CLIP thuần, CLIP + đồ thị và CLIP + phương pháp 2024.
\begin{itemize}
\item \textbf{Kênh đếm cạnh cho mức tăng lớn nhất.} Kênh này không có tham số huấn luyện và nâng R@1 val từ 24,33 lên 32,40.
Trên test, dòng chính hơn CLIP 4,93 điểm R@1 và 5,28 điểm MRR.
\item \textbf{NativE dùng riêng.} Điểm bộ ba của NativE nâng R@1 val 3,92 điểm so với CLIP; cosine embedding nâng 2,32 điểm.
Thứ tự này giữ nguyên trên test và ở cả 3 seed: kênh dùng quan hệ (\texttt{kgc}) cho kết quả cao hơn kênh chỉ so vector thực
thể (\texttt{emb}), nhưng cả hai thấp hơn kênh đếm cạnh. Như vậy, trong dòng chính của cấp~2, NativE không đóng góp; kết quả
của NativE được báo cáo ở phần tái lập (Bảng~\ref{tab:l2-repro}) và ở hai dòng dùng riêng.
\end{itemize}

\textbf{Độ nhạy theo $\alpha$} (Hình~\ref{fig:alpha}). Trên MKG-W, R@1 val giữ ở 32,40 trong cả đoạn $\alpha \in [0; 0{,}35]$
rồi giảm dần về 24,33 ở $\alpha = 1$; giá trị được chọn $\alpha = 0$ nằm trong vùng phẳng này. Khác với VG (đỉnh ở
$\alpha = 0{,}6$), trên MKG-W chỉ dùng đồ thị (32,40) cho kết quả cao hơn chỉ dùng CLIP (24,33). Với $\alpha = 0$, điểm cuối
chỉ còn số cạnh; các ứng viên cùng số cạnh giữ thứ tự của CLIP vì phép sắp xếp ổn định.

\textbf{Theo nhóm truy vấn.} Truy vấn không khớp được thực thể nào (189 câu val, 170 câu test) giữ nguyên kết quả của CLIP
theo thiết kế. Toàn bộ mức tăng nằm ở nhóm khớp được: R@1 val 20,91 $\to$ 31,15, test 22,71 $\to$ 28,81.

\textbf{Thứ tự kênh kgc được thêm.} Lần chạy val đầu tiên chỉ có kênh cosine, và val chọn $\beta = 0$ cho kênh này. Vì kênh
cosine không dùng quan hệ, kênh \texttt{kgc} được thêm để dùng hàm chấm có quan hệ của chính NativE. Định nghĩa kênh và dòng
chính mới được chốt trước khi chạy lại val và trước lần chấm test ngày 17/11. Lần chạy lại val giữ nguyên từng byte các file
kết quả cũ (\file{experiments/level2/val/README.md}).

\textbf{Các lần đọc test khác và một lần chạy lại độc lập.} (1) Khi kiểm tra code đánh giá, một lần chạy thử
\texttt{--split test --hand} với embedding giả đã in ra số test của dòng đếm cạnh, kể cả trên bộ câu tự viết. Dòng này không
có tham số nào ngoài $\alpha$, và $\alpha$ đã chốt trên val; không có gì được chỉnh sau đó. (2) Ngày 18/11, nhóm chạy lại độc
lập toàn bộ cấp~2 trên molab: tự huấn luyện NativE 3 seed (MRR rmask \pmstd{36,89}{0,27}, bản gốc 37,06), dùng đặc trưng
CLIP gốc thay cho bản nhóm trích lại, tự chọn $\alpha,\beta$ trên val rồi chấm test một lần. Val lại chọn $\alpha = 0$,
$\beta = 0$ ở cả 3 seed; dòng chính trên test cho đúng R@1 29,82, cosine NativE \pmstd{26,05}{0,51} và cạnh xáo trộn
\pmstd{21,41}{0,11}, gần với Bảng~\ref{tab:l2-results} và~\ref{tab:l2-ablation}. CLIP thuần trên test là 25,00 thay vì 24,89
vì khác bản đặc trưng. Kết quả và log của lần chạy này ở \file{experiments/level2/molab-rerun/}.

\subsubsection{Bộ 25 truy vấn tự viết: trực tiếp và gián tiếp}\label{sec:l2-hand}
\textbf{Cách xây và kiểm tra.} Bộ câu nằm ở \file{data/queries/mkgw_queries_hand.json}. Quy tắc được kiểm tra tự động bởi
\file{scripts/build_mkgw_hand_queries.py}; script từ chối ghi file nếu có câu vi phạm.
\begin{itemize}
\item \textbf{Trực tiếp} (10 câu): tả thực thể đáp án là gì, trông ra sao, không nêu tên nó và không dựa vào thực thể có tên
khác. Mỗi câu có đúng một đáp án trong kho ảnh test. Script kiểm tra mọi thực thể khớp được trong câu đều không có cạnh tới
đáp án, nên câu trực tiếp không có đường tắt qua đồ thị.
\item \textbf{Gián tiếp} (15 câu): nêu tên một thực thể khác và cần một quan hệ để tới đáp án, ví dụ ``a rock guitarist who
was a member of Pink Floyd''. Đáp án không gõ tay mà được tính từ mẫu (quan hệ, thực thể) của câu: mọi thực thể trong kho
ảnh test có đủ các cạnh đó, từ 1 tới 5 đáp án.
\item Tên đáp án không được xuất hiện trong câu. Đáp án tính trên toàn bộ triple MKG-W, vì đó là đáp án thật chứ không phải
đầu vào của mô hình. Kênh đồ thị lúc chấm vẫn chỉ thấy cạnh KGC-train. Bộ câu được commit một lần (16/11) và không sửa sau đó.
\end{itemize}

\begin{table}[t]
\centering
\caption{Bộ 25 truy vấn tự viết, ảnh của tập test làm ứng viên, dùng lại $\alpha,\beta$ của val.}
\label{tab:l2-hand}
\footnotesize
\setlength{\tabcolsep}{3pt}
\begin{tabular}{llcccc}
\toprule
Nhóm & Dòng & R@1 & R@5 & R@10 & MRR \\
\midrule
Trực tiếp (10) & CLIP & 80,00 & 90,00 & 90,00 & 85,00 \\
 & + đếm cạnh & 80,00 & 90,00 & 90,00 & 85,00 \\
Gián tiếp (15) & CLIP & 40,00 & 60,00 & 60,00 & 49,28 \\
 & + đếm cạnh & \textbf{80,00} & \textbf{86,67} & \textbf{86,67} & \textbf{83,33} \\
\bottomrule
\end{tabular}
\end{table}

\textbf{Kết quả} (Bảng~\ref{tab:l2-hand}). Với câu gián tiếp, kênh đếm cạnh nâng R@1 từ 40 lên 80 (6/15 $\to$ 12/15). Câu
trực tiếp không đổi: CLIP đã đúng 8/10, và theo cách xây, kênh đồ thị không cộng điểm cho ảnh đúng của các câu này. Nhóm tra
xem cạnh vàng của 15 câu gián tiếp nằm ở phần nào của KGC. 12/15 câu có ít nhất một đáp án nối với thực thể trong câu bằng
cạnh KGC-train, nên kênh đếm cạnh thấy được. 3/15 câu chỉ có cạnh vàng trong KGC-test, nên kênh này không thể giúp; ví dụ
``the rock band whose lead singer is Hayley Williams'' (đáp án Paramore). Do bộ câu chỉ có 25 câu, kết quả này mang tính
minh hoạ và không thay cho kết quả trên 892 truy vấn.

\textbf{Ví dụ lời giải thích.} Với câu ``the main railway station of Bordeaux'', \texttt{EntityLinker} khớp
\emph{Bordeaux}. Ảnh đúng là ga Bordeaux-Saint-Jean, có cạnh KGC-train P131 (\emph{located in the administrative territorial
entity}) tới Bordeaux. Lời giải thích của ứng viên là cạnh này, cũng là cạnh được đếm vào điểm đồ thị. Với câu về Pink
Floyd, có hai cạnh KGC-train được đếm: P463 (\emph{member of}, David Gilmour $\to$ Pink Floyd) và P527 (\emph{has part},
Pink Floyd $\to$ David Gilmour), nên điểm đồ thị của ảnh đúng là 2. Lần chấm ngày 17/11 chỉ lưu số tổng theo nhóm câu, không
lưu hạng của ảnh đúng theo từng câu, nên báo cáo không nêu hạng của từng câu.
